% Part: first-order-logic
% Chapter: natural-deduction
% Section: proving-things

\documentclass[../../../include/open-logic-section]{subfiles}

\begin{document}

\iftag{FOL}
      {\olfileid{fol}{ntd}{pro}}
      {\olfileid{pl}{ntd}{pro}}

\olsection{Voorbeelden van \usetoken{P}{derivation}}

\begin{ex}
We maken !!a{derivation} van de !!{sentence}
$(!A \land !B) \lif !A$.

Eerst schrijven we de gewenste conclusie onderaan de !!{derivation}.
\begin{prooftree}
\AxiomC{}
\UnaryInfC{$(!A\land !B) \lif !A$}
\end{prooftree}

Vervolgens zoeken we uit welke regel !!a{sentence} van deze vorm kan
opleveren. De !!{main operator}, dus het buitenste logische teken, van
de conclusie is $\lif$. Daarom proberen we met de regel \Intro{\lif}
bij de conclusie uit te komen. Schrijf tijdens het werk de gebruikte
aannamen op en zet de naam bij elke afleidingsregel. Dan kun je aan
het einde van het bewijs gemakkelijk controleren of alle aannamen zijn
!!{discharged}.
\begin{prooftree}
\AxiomC{$\Discharge{!A \land !B}{1}$}
\DeduceC{$!A$}
\DischargeRule{\Intro{\lif}}{1} 
\UnaryInfC{$(!A\land !B) \lif !A$}
\end{prooftree}

Nu vullen we de stappen van de aanname $!A \land !B$ naar $!A$ in.
Er is maar één logisch verbindingsteken waarmee we rekening moeten
houden: $\land$. Daarom gebruiken we de $\land$-eliminatieregel. Dat
geeft het volgende bewijs:
\begin{prooftree}
\AxiomC{$\Discharge{!A \land !B}{1}$}
\RightLabel{\Elim{\land}}
\UnaryInfC{$!A$}
\DischargeRule{\Intro{\lif}}{1} 
\UnaryInfC{$(!A\land !B) \lif !A$}
\end{prooftree}
We hebben nu een correcte !!{derivation} van $(!A \land
!B) \lif !A$.
\end{ex}

\begin{ex}
Nu maken we !!a{derivation} van $(\lnot !A \lor !B)
\lif (!A \lif !B)$.

Eerst schrijven we de gewenste conclusie onderaan de !!{derivation}.
\begin{prooftree}
\AxiomC{}
\UnaryInfC{$(\lnot !A \lor !B) \lif (!A \lif !B)$}
\end{prooftree}
We zoeken een logische regel die deze conclusie kan opleveren. Daarvoor
kijken we naar de logische verbindingstekens in de conclusie:
$\lnot$, $\lor$ en $\lif$. Op dit moment telt alleen het eerste
voorkomen van $\lif$. Dat is de !!{main operator}, dus het buitenste
logische teken, van de !!{sentence} in het eindsequent. $\lnot$,
$\lor$ en het tweede voorkomen van $\lif$ staan binnen het bereik van
een ander verbindingsteken. Die behandelen we later. We beginnen dus
met de regel \Intro{\lif}. Een correcte toepassing ziet er zo uit:
\begin{prooftree}
\AxiomC{$\Discharge{\lnot !A \lor !B}{1}$}
\DeduceC{$!A \lif !B$}
\DischargeRule{\Intro{\lif}}{1}
\UnaryInfC{$(\lnot !A \lor !B) \lif (!A \lif !B)$}
\end{prooftree}
We kunnen nu op twee manieren verdergaan. We kunnen van onder naar
boven blijven werken en de regel \Intro{\lif} nog eens toepassen. We
kunnen ook van boven naar beneden werken en een regel \Elim{\lor}
toepassen. We kiezen de tweede manier. De aanname
$\lnot !A \lor !B$ wordt de meest linkse uitgangsformule van
\Elim{\lor}. Voor een geldige toepassing van \Elim{\lor} moeten de
twee andere uitgangsformules hetzelfde zijn als de conclusie
$!A \lif !B$. We mogen elk daarvan afleiden uit een andere aanname:
uit een van de twee disjuncten van $\lnot !A \lor !B$. Onze
!!{derivation} ziet er dus zo uit:
\begin{prooftree}
\AxiomC{$\Discharge{\lnot !A \lor !B}{1}$}
\AxiomC{$\Discharge{\lnot !A}{2}$}
\DeduceC{$!A \lif !B$}
\AxiomC{$\Discharge{!B}{2}$}
\DeduceC{$!A \lif !B$}
\DischargeRule{\Elim{\lor}}{2}
\TrinaryInfC{$!A \lif !B$}
\DischargeRule{\Intro{\lif}}{1} 
\UnaryInfC{$(\lnot !A \lor !B) \lif (!A \lif !B)$}
\end{prooftree}

In elk van de twee takken rechts willen we $!A \lif !B$ !!{derive}.
Dat gaat het gemakkelijkst met \Intro{\lif}.
\begin{prooftree}
\AxiomC{$\Discharge{\lnot !A \lor !B}{1}$}
\AxiomC{$\Discharge{\lnot !A}{2}, \Discharge{!A}{3}$}
\DeduceC{$!B$}
\DischargeRule{\Intro{\lif}}{3}
\UnaryInfC{$!A \lif !B$}
\AxiomC{$\Discharge{!B}{2}, \Discharge{!A}{4}$}
\DeduceC{$!B$}
\DischargeRule{\Intro{\lif}}{4}
\UnaryInfC{$!A \lif !B$}
\DischargeRule{\Elim{\lor}}{2}
\TrinaryInfC{$!A \lif !B$}
\DischargeRule{\Intro{\lif}}{1} 
\UnaryInfC{$(\lnot !A \lor !B) \lif (!A \lif !B)$}
\end{prooftree}

Er ontbreken nog twee delen van de !!{derivation}. In het midden
hebben we !!{derivation}s van $!B$ uit $\lnot !A$ en $!A$ nodig.
Links hebben we er een uit $!A$ en $!B$ nodig. We beginnen
met het eerste geval. $\lnot !A$ en $!A$ zijn de twee
uitgangsformules van \Elim{\lnot}:
\begin{prooftree}
\AxiomC{$\Discharge{\lnot !A}{2}$}
\AxiomC{$\Discharge{!A}{3}$}
\RightLabel{\Elim{\lnot}}
\BinaryInfC{$\lfalse$}
\DeduceC{$!B$}
\end{prooftree}
Met \FalseInt krijgen we $!B$ als conclusie en kunnen we de tak
afmaken.
\begin{prooftree}
\AxiomC{$\Discharge{\lnot !A \lor !B}{1}$}
\AxiomC{$\Discharge{\lnot !A}{2}$}
\AxiomC{$\Discharge{!A}{3}$}
\RightLabel{\Intro{\lfalse}}
\BinaryInfC{$\lfalse$}
\RightLabel{\FalseInt}
\UnaryInfC{$!B$}
\DischargeRule{\Intro{\lif}}{3}
\UnaryInfC{$!A \lif !B$}
\AxiomC{$\Discharge{!B}{2}, \Discharge{!A}{4}$}
\DeduceC{$!B$}
\DischargeRule{\Intro{\lif}}{4}
\UnaryInfC{$!A \lif !B$}
\DischargeRule{\Elim{\lor}}{2}
\TrinaryInfC{$!A \lif !B$}
\DischargeRule{\Intro{\lif}}{1} 
\UnaryInfC{$(\lnot !A \lor !B) \lif (!A \lif !B)$}
\end{prooftree}

Kijk nu naar de meest rechtse tak. De definitie van !!{derivation}
\emph{staat toe dat we aannamen opheffen}, maar \emph{verplicht ons
daar niet toe}. We mogen $!B$ dus uit een van de aannamen $!A$ en
$!B$ afleiden zonder de andere te gebruiken. $!B$ uit~$!B$
!!{derive} is heel eenvoudig: $!B$ op zichzelf is al zo'n
!!{derivation}. Er is geen inferentie nodig. We kunnen de
aanname~$!A$ daarom gewoon weglaten.
\begin{prooftree}
\AxiomC{$\Discharge{\lnot !A \lor !B}{1}$}
\AxiomC{$\Discharge{\lnot !A}{2}$}
\AxiomC{$\Discharge{!A}{3}$}
\RightLabel{\Elim{\lnot}}
\BinaryInfC{$\lfalse$}
\RightLabel{\FalseInt}
\UnaryInfC{$!B$}
\DischargeRule{\Intro{\lif}}{3}
\UnaryInfC{$!A \lif !B$}
\AxiomC{$\Discharge{!B}{2}$}
\RightLabel{\Intro{\lif}}
\UnaryInfC{$!A \lif !B$}
\DischargeRule{\Elim{\lor}}{2}
\TrinaryInfC{$!A \lif !B$}
\DischargeRule{\Intro{\lif}}{1} 
\UnaryInfC{$(\lnot !A \lor !B) \lif (!A \lif !B)$}
\end{prooftree}
Let erop dat de meest rechtse \Intro{\lif}-inferentie in de voltooide
!!{derivation} geen enkele aanname opheft.
\end{ex}

\begin{ex}
Tot nu toe hadden we de regel \FalseCl{} niet nodig. Deze regel is
bijzonder. We mogen er namelijk een aanname mee opheffen die geen
formuledeel (sub-!!{formula}) van de conclusie van de regel is. De
regel hangt nauw samen met \FalseInt{}. \FalseInt{} is zelfs een
bijzonder geval van \FalseCl{}. Er bestaat een logica die
‘intuïtionistische logica’ heet. Daarin mag alleen \FalseInt{}
worden gebruikt. De regel \FalseCl{} is een laatste
uitweg als niets anders werkt. Stel bijvoorbeeld dat we
$!A \lor \lnot !A$ willen !!{derive}. Normaal zouden we proberen
$!A \lor \lnot !A$ met $\Intro{\lor}$ te !!{derive}. Dan zouden we
echter zonder aannamen
hetzij $!A$, hetzij $\lnot !A$ moeten !!{derive}. Dat kan niet.
\FalseCl{} biedt uitkomst.
\begin{prooftree}
  \AxiomC{$\Discharge{\lnot(!A \lor \lnot !A)}{1}$}
  \DeduceC{$\lfalse$}
  \DischargeRule{\FalseCl}{1}
  \UnaryInfC{$!A \lor \lnot !A$}
\end{prooftree}
Nu zoeken we !!a{derivation} van $\lfalse$ uit $\lnot(!A \lor
\lnot !A)$. $\lfalse$ is de conclusie van $\Elim{\lnot}$, dus we
kunnen die regel proberen:
\begin{prooftree}
  \AxiomC{$\Discharge{\lnot(!A \lor \lnot !A)}{1}$}
  \DeduceC{$\lnot !A$}
  \AxiomC{$\Discharge{\lnot(!A \lor \lnot !A)}{1}$}
  \DeduceC{$!A$}
  \RightLabel{\Elim{\lnot}}
  \BinaryInfC{$\lfalse$}
  \DischargeRule{\FalseCl}{1}
  \UnaryInfC{$!A \lor \lnot !A$}
\end{prooftree}
Om !!a{derivation} van~$\lnot !A$ te vinden, passen we
~$\Intro{\lnot}$ toe:
\begin{prooftree}
  \AxiomC{$\Discharge{\lnot(!A \lor \lnot !A)}{1}, \Discharge{!A}{2}$}
  \DeduceC{$\lfalse$}
  \DischargeRule{\Intro{\lnot}}{2}
  \UnaryInfC{$\lnot !A$}
  \AxiomC{$\Discharge{\lnot(!A \lor \lnot !A)}{1}$}
  \DeduceC{$!A$}
  \RightLabel{\Elim{\lnot}}
  \BinaryInfC{$\lfalse$}
  \DischargeRule{\FalseCl}{1}
  \UnaryInfC{$!A \lor \lnot !A$}
\end{prooftree}
Hier krijgen we $\lfalse$ gemakkelijk met $\Elim{\lnot}$. We passen
die regel toe op de aanname $\lnot(!A \lor \lnot !A)$ en op
$!A \lor \lnot !A$. Deze tweede formule volgt uit onze nieuwe aanname
$!A$ met~$\Intro{\lor}$:
\begin{prooftree}
  \AxiomC{$\Discharge{\lnot(!A \lor \lnot !A)}{1}$}
  \AxiomC{$\Discharge{!A}{2}$}
  \RightLabel{\Intro{\lor}}
  \UnaryInfC{$!A \lor \lnot !A$}
  \RightLabel{\Elim{\lnot}}
  \BinaryInfC{$\lfalse$}
  \DischargeRule{\Intro{\lnot}}{2}
  \UnaryInfC{$\lnot !A$}
  \AxiomC{$\Discharge{\lnot(!A \lor \lnot !A)}{1}$}
  \DeduceC{$!A$}
  \RightLabel{\Elim{\lnot}}
  \BinaryInfC{$\lfalse$}
  \DischargeRule{\FalseCl}{1}
  \UnaryInfC{$!A \lor \lnot !A$}
\end{prooftree}
Rechts volgen we dezelfde stappen, maar krijgen we $!A$ met
~\FalseCl:
\begin{prooftree}
  \AxiomC{$\Discharge{\lnot(!A \lor \lnot !A)}{1}$}
  \AxiomC{$\Discharge{!A}{2}$}
  \RightLabel{\Intro{\lor}}
  \UnaryInfC{$!A \lor \lnot !A$}
  \RightLabel{\Elim{\lnot}}
  \BinaryInfC{$\lfalse$}
  \DischargeRule{\Intro{\lnot}}{2}
  \UnaryInfC{$\lnot !A$}
  \AxiomC{$\Discharge{\lnot(!A \lor \lnot !A)}{1}$}
  \AxiomC{$\Discharge{\lnot !A}{3}$}
  \RightLabel{\Intro{\lor}}
  \UnaryInfC{$!A \lor \lnot !A$}
  \RightLabel{\Elim{\lnot}}
  \BinaryInfC{$\lfalse$}
  \DischargeRule{\FalseCl}{3}
  \UnaryInfC{$!A$}
  \RightLabel{\Elim{\lnot}}
  \BinaryInfC{$\lfalse$}
  \DischargeRule{\FalseCl}{1}
  \UnaryInfC{$!A \lor \lnot !A$}
\end{prooftree}
\end{ex}

\begin{prob}
Geef !!{derivation}s die het volgende laten zien:
\begin{enumerate}
\item $!A \land (!B \land !C) \Proves (!A \land !B) \land !C$.
\item $!A \lor (!B \lor !C) \Proves (!A \lor !B) \lor !C$.
\item $!A \lif (!B \lif !C) \Proves !B \lif (!A \lif !C)$.
\item $!A \Proves \lnot\lnot !A$.
\end{enumerate}
\end{prob}

\begin{prob}
Geef !!{derivation}s die het volgende laten zien:
\begin{enumerate}
\item $(!A \lor !B) \lif !C \Proves !A \lif !C$.
\item $(!A \lif !C) \land (!B \lif !C) \Proves (!A \lor !B) \lif !C$.
\item $\Proves \lnot(!A \land \lnot !A)$.
\item $!B \lif !A \Proves \lnot !A \lif \lnot !B$.
\item $\Proves (!A \lif \lnot !A) \lif \lnot !A$.
\item $\Proves \lnot(!A \lif !B) \lif \lnot !B$.
\item $!A \lif !C \Proves \lnot (!A \land \lnot !C)$.
\item $!A \land \lnot !C \Proves \lnot (!A \lif !C)$.
\item $!A \lor !B, \lnot !B \Proves !A$.
\item $\lnot !A \lor \lnot !B \Proves \lnot(!A \land !B)$.
\item $\Proves (\lnot !A \land \lnot !B) \lif\lnot(!A \lor !B)$.
\item $\Proves \lnot(!A \lor !B) \lif (\lnot !A \land \lnot !B)$.
\end{enumerate}
\end{prob}

\begin{prob}
Geef !!{derivation}s die het volgende laten zien:
\begin{enumerate}
\item $\lnot(!A \lif !B) \Proves !A$.
\item $\lnot(!A \land !B) \Proves \lnot !A \lor \lnot !B$.
\item $!A \lif !B \Proves \lnot !A \lor !B$.
\item $\Proves \lnot \lnot !A \lif !A$.
\item $!A \lif !B, \lnot !A \lif !B \Proves !B$.
\item $(!A \land !B) \lif !C \Proves (!A \lif !C) \lor (!B \lif !C)$.
\item $(!A \lif !B) \lif !A \Proves !A$.
\item $\Proves (!A \lif !B) \lor (!B \lif !C)$.
\end{enumerate}
(Hiervoor is steeds de regel $\FalseCl$ nodig.)
\end{prob}
\end{document}
